\documentclass{article}
\usepackage[utf8]{inputenc}
\usepackage[T1]{fontenc}
\usepackage{amsmath, amssymb}
\usepackage{bm}
\usepackage{hyperref, xcolor}
\usepackage[
backend=biber,
sorting=none,
style=numeric-comp
]{biblatex} %Imports biblatex package
\addbibresource{references.bib}

\renewcommand{\Re}{\operatorname{Re}}
\renewcommand{\Im}{\operatorname{Im}}
\newcommand{\ettriangle}{ET$-\Delta$}
\newcommand{\etaligned}{ET-2L-aligned}
\newcommand{\etmisaligned}{ET-2L-misaligned}
\newcommand{\fref}{f_{\rm{ref}}}
\newcommand{\phiref}{\phi_{\rm{ref}}}

% For comments
\newcommand{\bernardo}[1]{{\color{red} [BV: #1]}}
\newcommand{\rs}[1]{{\color{blue} [RS: #1]}}

\title{SGWB}
\author{Bernardo Veronese}

\begin{document}

\maketitle

\section{Introduction}

\section{Theory}

We label our hyperparameter vector by $\Lambda$, comprising astrophysical
parameters $\Lambda_a$ and cosmological parameters $\Lambda_c$. The likelihood
is given by the squared SNR of the residuals:

\begin{equation}
	\mathcal{L} \propto 2T \sum_{a,b}\int_{f_{\rm{min}}}^{f_{\rm{max}}} \left(S_h(\Lambda) - S_{h,\rm{data}}\right)^2 \frac{\Gamma^2_{ab}(f)}{S_{n,a}(f)S_{n,b}(f)} df.
	\label{eq:likelihood}
\end{equation}

The spectral density of the data, $S_{h,\rm{data}}$, is the sum of the
contributions of all astrophysical sources, see~\cite{Belgacem:2024ohp}:

\begin{equation}
	S_h^{\rm{astro}}(f) = \frac{1}{T} \sum_{i=1}^{N} |\tilde{h}_{i,+}(f, \theta_i)|^2 + |\tilde{h}_{i,\times}(f, \theta_i)|^2,
	\label{eq:data-spectral-density}
\end{equation}

where $\theta_i$ denotes the intrinsic and extrinsic parameters of the $i$-th
source, and $N$ is the total number of sources.

In practice, the sum will be carried out over a known catalog, like those
provided by the COBA
paper~\cite{Branchesi:2023mws}\footnote{\url{https://apps.et-gw.eu/tds/?r=18321}}.
The averages over parameters are then replaced by the average over the catalog,
which is a Monte Carlo sampling of the underlying population.

In general, the theoretical expression for $S_h^{\rm{astro}}$ is given by:

\begin{equation}
	\begin{aligned}
		S_h^{\rm{astro}}(f; \Lambda_a, \Lambda_c) & = \left \langle \left[|\tilde{h}_{+}(f, \theta)|^2 + |\tilde{h}_{\times}(f, \theta)|^2\right] \right \rangle,                                                       \\
		                                          & = \int dz d\theta_i \, p(z, \theta_i | \Lambda) \left \langle \left[|\tilde{h}_{+}(f, \theta)|^2 + |\tilde{h}_{\times}(f, \theta)|^2\right] \right \rangle_{\iota},
	\end{aligned}
	\label{eq:theoretical-spectral-density}
\end{equation}

The conditional probability $p(z, \theta_i | \Lambda)$ is the probability of
having a source with parameters $\theta_i$ at redshift $z$, given the
hyperparameters $\Lambda$. We assume that the distribution of intrinsic
parameters is independent of redshift, so that we can write $p(z, \theta_i |
	\Lambda) = p(z | \Lambda) p(\theta_i | \Lambda)$. The redshift distribution is
given by:

\begin{equation}
	p(z | \Lambda) = \frac{R(z; \Lambda)}{1+z} \frac{dV_c}{dz},
	\label{eq:redshift-distribution}
\end{equation}

where $R(z; \Lambda)$ is the merger rate per comoving volumeat redshift $z$. We
model it with the functional form inspired by the star formation rate density,

\begin{equation}
	R(z; \Lambda) = R_0 \frac{(1+z)^{\alpha}}{1 + \left(\frac{1+z}{1+z_p}\right)^{\alpha+\beta}},
	\label{eq:merger-rate}
\end{equation}

where $\alpha, \beta, z_p \in \Lambda_a$, and $R_0$ is the local merger rate
density given in $\rm{Mpc}^{-3} \, \rm{yr}^{-1}$. The term $dV_c/dz$ is the
jacobian for the change of variable from comoving volume to redshift. The $1+z$
factor accounts for the time dilation between the source frame and the detector
frame.

\section{TODOs}
\begin{itemize}
	\item Read~\cite{Capurri:2023jsy} thoroughly
	\item Read~\cite{Belgacem:2018lbp} thoroughly
	\item Read ET Blue Book section 2.3.2
	\item Learn about PLS (power law integrated sensitivty curve)
	\item Download BNS catalog from COBA paper
	\item Think about how to implement SGWB-related stuff (talk to Nicolo about it)
	\item Run benchmark with~\cite{Capurri:2023jsy}
\end{itemize}

\printbibliography

/appendix

\section{Marginalization over the inclination angle}

On the assumption that the waveform is dominate by the $(2, 2)$ mode, the
marginalization over the inclination angle $\iota$ can be performed
analytically. We write the waveform polarizations as

\begin{equation}
	\begin{aligned}
		h_+(f, \theta)              & = A(f, \theta) \frac{1 + \cos^2\iota}{2}, \\
		\quad h_{\times}(f, \theta) & = A(f, \theta) \cos\iota,
	\end{aligned}
	\label{eq:waveform-polarizations-inclination}
\end{equation}

Let $I(f, \theta) = |\tilde{h}_{+}(f, \theta)|^2 + |\tilde{h}_{\times}(f,
	\theta)|^2$. We can factor out the dependence on the inclination angle as
follows:

The average over the inclination angle is given by the integral

\begin{equation}
	\langle I(f, \theta) \rangle_\iota =	\frac{1}{2} \int_0^\pi d\iota p(\iota) \left[|\tilde{h}_{+}(f, \theta)|^2 + |\tilde{h}_{\times}(f, \theta)|^2\right].
	\label{eq:inclination-averaged-waveform}
\end{equation}

The natural choice of prior for the inclination angle is a uniform distribution
in $\cos \iota$, which corresponds to an isotropic distribution of sources in
the sky. Thus, Eq.~\eqref{eq:inclination-averaged-waveform} becomes:

\begin{align}
	\nonumber
	\langle I(f, \theta) \rangle_\iota & = \frac{1}{2} A(f, \theta) \int d\cos\iota \frac{1}{2} \left[ \left(\frac{1 + \cos^2\iota}{2}\right)^2 + \cos^2\iota \right] \\
	\nonumber
	                                   & = \frac{4}{5} A(f, \theta)                                                                                                   \\
	                                   & = \frac{2}{5} I(f, \theta, \iota = 0).
	\label{eq:}
\end{align}
\end{document}
